% ARCHITECTURE.tex — how the WebCV site works, and why.
% Companion docs: SPEC.md (requirements) and PLAN.md (execution tracker).
% Compile with: tectonic ARCHITECTURE.tex   (or pdflatex/xelatex, run twice
% for the table of contents and cross-references to resolve).

\documentclass[11pt]{article}

% ---------------------------------------------------------------------------
% Packages
% ---------------------------------------------------------------------------
\usepackage[a4paper,margin=2.4cm]{geometry}
\usepackage[utf8]{inputenc}
\usepackage[T1]{fontenc}
\usepackage{titlesec}
\usepackage{enumitem}
\usepackage{longtable}
\usepackage{array}
\usepackage{booktabs}
\usepackage{xcolor}
\usepackage{tcolorbox}
\usepackage{fancyhdr}
\usepackage{parskip}
\usepackage{tikz}
\usepackage{listings}
\usepackage{hyperref}
\usetikzlibrary{arrows.meta, positioning, fit, backgrounds}
\tcbuselibrary{skins,breakable}

% ---------------------------------------------------------------------------
% Palette — matches the site's own three-anime blend (SPEC.md §9.1):
% Naruto orange (primary), One Piece straw-hat gold, Food Wars crimson.
% ---------------------------------------------------------------------------
\definecolor{narutoorange}{HTML}{E8611C}
\definecolor{deeporange}{HTML}{C74C10}
\definecolor{strawgold}{HTML}{8A6116}
\definecolor{shokugekicrimson}{HTML}{C81E3A}
\definecolor{inkgray}{HTML}{232323}
\definecolor{mutedgray}{HTML}{5B5B5B}
\definecolor{bgcream}{HTML}{FBF6EE}
\definecolor{bordergray}{HTML}{E7DED0}
\definecolor{codebg}{HTML}{F5F1E6}

\hypersetup{
  colorlinks=true,
  linkcolor=deeporange,
  urlcolor=narutoorange,
  citecolor=strawgold,
  pdftitle={ARCHITECTURE --- WebCV},
  pdfauthor={Sudarshana Chaitanya B N},
  pdfsubject={System architecture and design decisions for the WebCV site},
}

% Section styling
\titleformat{\section}
  {\normalfont\Large\bfseries\color{narutoorange}}
  {\thesection}{0.6em}{}
\titleformat{\subsection}
  {\normalfont\large\bfseries\color{deeporange}}
  {\thesubsection}{0.6em}{}
\titleformat{\subsubsection}
  {\normalfont\bfseries\color{inkgray}}
  {\thesubsubsection}{0.6em}{}

\pagestyle{fancy}
\fancyhf{}
\renewcommand{\headrulewidth}{0.4pt}
\renewcommand{\footrulewidth}{0pt}
\fancyhead[L]{\small\color{mutedgray} ARCHITECTURE --- WebCV}
\fancyhead[R]{\small\color{mutedgray} \thepage}

\lstset{
  basicstyle=\ttfamily\small,
  backgroundcolor=\color{codebg},
  frame=single,
  rulecolor=\color{bordergray},
  breaklines=true,
  columns=fullflexible,
  keepspaces=true,
  xleftmargin=4pt,
  xrightmargin=4pt,
}

% A callout box for "why it matters" asides
\newtcolorbox{whybox}[1][]{
  enhanced, breakable,
  colback=bgcream, colframe=narutoorange,
  boxrule=0.6pt, arc=2pt,
  left=6pt, right=6pt, top=4pt, bottom=4pt,
  title={\textbf{Why it matters}},
  fonttitle=\small\bfseries\color{deeporange},
  coltitle=deeporange,
  attach boxed title to top left={xshift=6pt,yshift=-2pt},
  boxed title style={colback=bgcream, boxrule=0pt},
}

% \file{a/b/c.ext} — typesets a path in monospace, with a break opportunity
% after every "/" so long paths wrap inside narrow table columns instead of
% overflowing the margin. Plain \allowbreak (no catcode tricks), so it's
% safe to use even inside moving arguments like \subsection{}.
\makeatletter
\newcommand{\file}[1]{\texttt{\color{shokugekicrimson}\@filebreak#1/\@stop}}
\def\@filebreak#1/#2\@stop{%
  #1%
  \def\@tmp{#2}%
  \ifx\@tmp\@empty\else/\allowbreak\@filebreak#2\@stop\fi
}
\makeatother

\title{
  \vspace{-1.5cm}
  {\Large\bfseries\color{narutoorange} ARCHITECTURE}\\[4pt]
  {\large How the WebCV site works, and why}
}
\author{Sudarshana Chaitanya B N}
\date{Last updated: \today}

% ===========================================================================
\begin{document}

\maketitle
\thispagestyle{fancy}

\begin{whybox}[title={\textbf{What this document is for}}]
This explains the WebCV project well enough to walk someone else --- a
teammate, a reviewer, an interviewer --- through it without relying on
memory. It covers what each piece does, why it was built that way instead
of some other way, and what actually broke along the way. \file{SPEC.md}
is the requirements doc (what the site must do); \file{PLAN.md} is the
execution tracker (what's done, what's open); this document is the
``how it's built, and why.''
\end{whybox}

\tableofcontents

% ===========================================================================
\section{The 30-Second Pitch}
\label{sec:pitch}

\begin{quote}
\itshape
``It's a personal portfolio site built with Astro, a static site
generator. Content --- CV, work experience, and projects --- lives as
plain YAML/Markdown files, completely separate from the page code. At
build time, Astro reads that data and pre-renders every page to plain
HTML/CSS, which gets deployed to GitHub Pages via a GitHub Actions
pipeline that runs automatically on every push. Because content and code
are separate, adding a new project or job later means editing one YAML
file --- no HTML, no redesign, no redeploy button to remember to press.''
\end{quote}

That one paragraph answers ``walk me through your project'' in one
breath, and every section below expands on a piece of it.

% ===========================================================================
\section{Big Picture: From a YAML File to a Live URL}
\label{sec:bigpicture}

\begin{center}
\resizebox{\textwidth}{!}{%
\begin{tikzpicture}[
  node distance=10mm and 8mm,
  box/.style={draw, rounded corners, thick, minimum height=1.1cm,
    align=center, font=\small, inner sep=6pt},
  data/.style={box, draw=strawgold, fill=strawgold!8},
  code/.style={box, draw=narutoorange, fill=narutoorange!8},
  outbox/.style={box, draw=shokugekicrimson, fill=shokugekicrimson!8},
  arrow/.style={-{Stealth[length=2.5mm]}, thick, draw=inkgray}
]
  \node[data] (yaml) {content/experience.yaml\\content/projects.yaml};
  \node[code, right=of yaml] (loader) {src/lib/content.ts\\\footnotesize (reads \& types YAML)};
  \node[code, right=of loader] (pages) {src/pages/*.astro\\\footnotesize (templates loop over data)};
  \node[outbox, right=of pages] (build) {astro build\\\footnotesize (static site generator)};
  \node[outbox, below=12mm of build] (dist) {dist/*.html + *.css};
  \node[code, left=of dist] (ci) {GitHub Actions\\\footnotesize (CI/CD pipeline)};
  \node[data, left=of ci] (pages_host) {GitHub Pages\\\footnotesize (static file host)};
  \node[outbox, left=of pages_host] (url) {dasa108.github.io\\/WebCV/};

  \draw[arrow] (yaml) -- (loader);
  \draw[arrow] (loader) -- (pages);
  \draw[arrow] (pages) -- (build);
  \draw[arrow] (build) -- (dist);
  \draw[arrow] (dist) -- (ci);
  \draw[arrow] (ci) -- (pages_host);
  \draw[arrow] (pages_host) -- (url);
\end{tikzpicture}%
}
\end{center}

Two things to internalize, because they come up in almost every
follow-up question:

\begin{enumerate}
  \item \textbf{Everything happens at build time, not request time.}
  There is no server running when someone visits the site.
  \texttt{astro build} runs once, in GitHub's CI runner --- not in the
  visitor's browser --- and produces plain HTML/CSS files, which is what
  gets hosted. This is \textbf{static site generation (SSG)}.
  \item \textbf{Content and code are two layers that only meet at build
  time.} \file{content/*.yaml} doesn't know anything about HTML.
  \file{src/pages/*.astro} doesn't know anything about \emph{this
  specific} internship or project --- it only knows how to render
  \emph{any} entry matching the shape defined in \file{src/lib/content.ts}.
  That separation is the whole reason adding a project is a one-file edit
  instead of a code change (\S\ref{sec:contentmodel}).
\end{enumerate}

% ===========================================================================
\section{Why Astro? What a Static Site Generator Is}
\label{sec:whyastro}

\textbf{The problem to solve:} almost all of this site's content changes
rarely (a CV doesn't change daily), and there's no login, no database, no
server-side logic --- it's fundamentally a document, not an application.
Serving a full client-rendered JavaScript app for that solves a problem
that doesn't exist here, at the cost of a slower first load and more
moving parts.

\subsection{The Three Real Options Weighed}

\begin{longtable}{@{}p{3.1cm}p{4.6cm}p{6.8cm}@{}}
\toprule
\textbf{Option} & \textbf{What it is} & \textbf{Why not chosen} \\
\midrule
\endhead
Plain HTML\slash CSS\slash JS, no build step &
  Hand-write every page &
  No templating --- adding a project means copy-pasting a whole card's
  HTML by hand, every time. Doesn't scale. \\[4pt]
React / Next.js &
  Full JS framework, client + server rendering &
  Overkill for a mostly-static content site; ships a JS framework
  runtime to the browser for pages that don't need interactivity. \\[4pt]
\textbf{Astro (chosen)} &
  Static site generator with component templating &
  Renders to plain HTML at build time (fast, simple hosting) but still
  gives reusable components and data loops --- templating power without
  shipping a framework runtime to the browser. \\
\bottomrule
\end{longtable}

\textbf{What ``static site generator'' means, concretely:} a program
that takes templates plus data and produces \texttt{.html} files
\emph{before} anyone visits the site, instead of generating HTML
on-the-fly per visitor (server-side rendering, SSR) or in the visitor's
browser with JavaScript (a single-page app). Astro's default mode ---
used here explicitly via \texttt{output: 'static'} in
\file{astro.config.mjs} --- is SSG: build once, serve static files
forever, until the next build.

\begin{whybox}
Static files are the simplest, cheapest, fastest thing to host --- a CDN
can cache them forever, there's no server to crash or patch, and GitHub
Pages hosts them for free. The tradeoff (rebuild to see new content) is a
non-issue because publishing a new build is automatic (\S\ref{sec:cicd})
--- it's just a \texttt{git push}.
\end{whybox}

% ===========================================================================
\section{The Content Model: Why \texttt{content/} Is Separate From \texttt{src/}}
\label{sec:contentmodel}

This is the single most important design decision in the project ---
\textbf{separation of concerns}: data kept separate from presentation.
It's a general software design principle, not an Astro-specific trick.

\textbf{The requirement:} adding a new internship or project later
should mean editing \emph{one file}, never touching page templates or
CSS.

\subsection{How It's Implemented}

\begin{itemize}[leftmargin=1.4em]
  \item \file{content/experience.yaml} and \file{content/projects.yaml}
  each hold a flat YAML \emph{list} of entries. Each entry is a plain
  object --- a role, an organization, bullet points, tags. Nothing in
  these files knows about HTML, colors, or layout.
  \item \file{src/lib/content.ts} is the \emph{only} code that touches
  these files. It reads them with Node's filesystem API and parses them
  with the \texttt{js-yaml} library, exposing typed functions:
  \texttt{getExperience()}, \texttt{getProjects()},
  \texttt{getProjectBySlug(slug)}.
  \item The page templates (\file{src/pages/*.astro}) call those
  functions and loop over whatever comes back --- they render ``however
  many entries exist,'' never ``these specific $N$ entries.''
\end{itemize}

So the flow for ``add a new project'' is: open \file{projects.yaml},
copy the \texttt{TEMPLATE} comment block at the top, fill in the fields,
paste it into the list, commit, push. The Projects page shows one more
card automatically, with zero code changes.

\subsection{Why YAML, Not a Database or a CMS?}

\begin{itemize}[leftmargin=1.4em]
  \item A database needs a server to run and a network call to query ---
  overkill for a few dozen KB of text that changes maybe once a month.
  \item A headless CMS (Contentful, Sanity) adds an external account, an
  API key, and a network dependency for content that's really just plain
  text files. YAML in the same git repo means the content is
  version-controlled, diffable, and needs no third-party service.
  \item YAML specifically, over JSON, because it supports comments ---
  which is what makes the in-file \texttt{TEMPLATE} block and
  field-by-field explanations possible. That self-documenting schema is
  what makes it usable by someone (future me) who doesn't remember the
  exact field names months later.
\end{itemize}

\subsection{Why Not Astro's Built-In ``Content Collections''?}

Content Collections expect one file per entry, living inside
\file{src/content/}. The existing ``one list file per type'' shape was
kept deliberately instead, for two reasons: (1) it's simpler to reason
about --- ``everything is in this one file'' vs. ``a new file per entry
plus a schema file'' --- and (2) the whole point of a scratch-built
loader (\file{content.ts}) is full control: no framework magic between
the YAML and what actually renders, which keeps debugging a matter of a
single function call you can log.

% ===========================================================================
\section{Repository Tour: What Every File Does, and Why It's There}
\label{sec:repotour}

The table below covers every path in the repository. The three-column
shape mirrors the question worth asking of any file in a project: what
is it, and why does it need to exist at all (as opposed to being deleted
or merged into something else)?

\begin{longtable}{@{}p{4.3cm}p{5.1cm}p{5.1cm}@{}}
\toprule
\textbf{Path} & \textbf{What it is} & \textbf{Why it's needed} \\
\midrule
\endfirsthead
\toprule
\textbf{Path} & \textbf{What it is} & \textbf{Why it's needed} \\
\midrule
\endhead
\file{SPEC.md} &
  Requirements doc: page structure, color palette, responsiveness/a11y
  targets, acceptance criteria. &
  Separates ``what must this do'' from ``how do I build it,'' and
  records decisions (breakpoints, palette) that would otherwise live
  only in memory and drift. \\[6pt]
\file{PLAN.md} &
  Execution tracker: what's done, what's open, the decision log, the
  playbook for adding content later. &
  \texttt{SPEC.md} barely changes; this file changes constantly and
  records \emph{why} a choice was made once, so it's not re-debated in
  every future conversation about the project. \\[6pt]
\file{ARCHITECTURE.md} / \file{ARCHITECTURE.tex} (this doc) &
  How the code works and why --- tech-stack decisions, code walkthrough,
  CI/CD, real bugs hit and fixed. &
  Lets the project be explained to someone else without re-deriving
  every decision from scratch. \\[6pt]
\file{README.md} &
  Entry point: what the project is, links to the other docs, how to run
  it locally. &
  The first thing anyone (including future me) reads --- has to orient a
  stranger in under a minute. \\[6pt]
\file{package.json} &
  Declares dependencies (\texttt{astro}, \texttt{js-yaml}) and the
  \texttt{npm run \{dev,build,preview\}} scripts. &
  The standard Node.js manifest --- without it, neither a human nor CI
  knows what to install or how to build. \\[6pt]
\file{package-lock.json} &
  Exact, pinned dependency version tree. &
  Makes builds reproducible: CI (\texttt{npm ci}) installs \emph{exactly}
  what was tested locally, not ``whatever satisfies the version range
  today.'' \\[6pt]
\file{astro.config.mjs} &
  Site config: \texttt{output: 'static'}, the GitHub Pages
  \texttt{site}/\texttt{base} URL. &
  Astro needs to know it's deploying to a subpath
  (\texttt{/WebCV/}) rather than a domain root, so every internal link
  and asset path resolves correctly. \\[6pt]
\file{tsconfig.json} &
  TypeScript compiler config, extending Astro's strict preset. &
  Turns on type-checking across \texttt{.ts}/\texttt{.astro} files ---
  catches a misspelled YAML field or wrong prop type before it ships. \\[6pt]
\file{.gitignore} &
  Excludes \texttt{dist/}, \texttt{.astro/}, \texttt{node\_modules/},
  \texttt{.env*} from version control. &
  Build output and dependencies are regenerable and often huge; they
  don't belong in git history. \\[6pt]
\file{.github/workflows/deploy.yml} &
  GitHub Actions CI/CD pipeline: build then deploy on every push to
  \texttt{main}. &
  Turns ``publish a change'' into ``git push'' --- see \S\ref{sec:cicd}. \\[6pt]
\file{content/experience.yaml} &
  Structured list of every work entry (internships, roles). &
  The actual data source for the Experience page --- see
  \S\ref{sec:contentmodel}. \\[6pt]
\file{content/projects.yaml} &
  Structured list of every project (title, summary, tags, links,
  \texttt{featured} flag). &
  The actual data source for the Projects grid and each project's detail
  page. \\[6pt]
\file{content/cv.md}, \file{internships.md}, \file{projects.md},
\file{extracurriculars.md} &
  Human-readable prose notes --- the raw material the YAML files were
  distilled from. &
  Kept as reference/source-of-truth prose; not parsed by the site
  directly, but useful when writing a new YAML entry from scratch. \\[6pt]
\file{content/*.pdf} &
  Original source documents (CV versions, project/internship reports). &
  Kept in \texttt{content/} deliberately \emph{not} web-accessible ---
  private working material, distinct from the specific PDFs chosen to be
  publicly linked (copied into \texttt{public/documents/} instead). \\[6pt]
\file{public/favicon.svg} &
  Browser-tab icon. &
  Small brand touch; SVG keeps it crisp at any size with no image
  processing needed. \\[6pt]
\file{public/resume.pdf} &
  The downloadable resume, linked from the hero and Contact section. &
  Recruiters expect a one-click PDF, not just a webpage, per
  \texttt{SPEC.md}'s acceptance criteria. \\[6pt]
\file{public/documents/*.pdf} &
  Report/slide PDFs linked from specific project entries. &
  Lives in \texttt{public/} (not \texttt{content/}) specifically because
  it \emph{must} be web-accessible --- see the \texttt{content/} vs.
  \texttt{public/} distinction below, and Bug 2 in \S\ref{sec:bugs}. \\[6pt]
\file{public/images/\{experience,projects\}/} &
  Org logos / project thumbnails, referenced by optional YAML fields. &
  Optional visual polish per entry, without forcing every entry to have
  one (fields are nullable). \\[6pt]
\file{public/images/profile.png} &
  Hero profile photo. &
  Referenced directly by \texttt{index.astro}; fails gracefully (hides
  itself, no broken-image icon) if absent, so the site never looks
  broken while this is still being decided. \\[6pt]
\file{src/lib/content.ts} &
  The data layer: reads/parses/types \texttt{content/*.yaml}. &
  The \emph{only} place YAML is touched --- single source of truth for
  ``what an experience/project entry looks like.'' See
  \S\ref{sec:codewalkthrough}. \\[6pt]
\file{src/styles/global.css} &
  Design tokens (CSS custom properties), dark mode, responsive base
  styles. &
  One place to define color/spacing \emph{by role}; every component
  reuses the same tokens instead of hardcoding values --- see
  \S\ref{sec:codewalkthrough}. \\[6pt]
\file{src/components/Layout.astro} &
  Page shell: \texttt{<html>/<head>} boilerplate, meta tags, Open Graph
  tags, nav, footer, \texttt{<slot />}. &
  Written once, so every page gets consistent SEO/meta handling instead
  of duplicating \texttt{<head>} content five times. \\[6pt]
\file{src/components/Nav.astro} &
  Site navigation, active-page highlighting. &
  Builds links from \texttt{import.meta.env.BASE\_URL} rather than a
  hardcoded path, so it keeps working if the deploy target changes. \\[6pt]
\file{src/components/Footer.astro} &
  Site footer. &
  Shared chrome, same reasoning as \texttt{Layout.astro}. \\[6pt]
\file{src/components/ExperienceCard.astro},
\file{ProjectCard.astro} &
  Render one experience/project entry as a card. &
  Take one typed entry as a prop and know nothing about \emph{how many}
  times they're used --- the page does the looping (\S\ref{sec:contentmodel}). \\[6pt]
\file{src/pages/index.astro} &
  The homepage: hero, about, skills, highlights, extracurriculars,
  contact. &
  Astro's file-based routing --- this file \emph{is} the route \texttt{/}. \\[6pt]
\file{src/pages/experience.astro} &
  The \texttt{/experience} page: loops over \texttt{getExperience()}. &
  One template, however many entries currently exist --- see
  \S\ref{sec:codewalkthrough}. \\[6pt]
\file{src/pages/projects/index.astro} &
  The \texttt{/projects} grid page. &
  Same pattern as Experience, applied to \texttt{getProjects()}. \\[6pt]
\file{src/pages/projects/[slug].astro} &
  Dynamic route: one generated page per project. &
  Uses \texttt{getStaticPaths()} so a new YAML entry becomes a new live
  page with zero route registration --- the direct mechanism behind
  ``one YAML entry = one live page.'' \\
\bottomrule
\end{longtable}

\subsection{\texttt{content/} vs. \texttt{public/}: Why Both Exist}

They look similar (both hold ``data'') but serve opposite purposes:

\begin{itemize}[leftmargin=1.4em]
  \item \textbf{\texttt{public/}} is copied \emph{verbatim} into the
  deployed site --- anything here gets its own URL
  (\file{public/resume.pdf} $\to$ \file{/WebCV/resume.pdf}).
  \item \textbf{\texttt{content/}} is \emph{never} copied to the
  deployed site --- it's read at build time by \texttt{content.ts},
  turned into HTML, and the original files stay private to the repo.
\end{itemize}

This is deliberate: raw PDFs and prose notes in \texttt{content/} don't
need to be individually web-accessible; only the specific documents
chosen to link are (copied into \texttt{public/documents/}). Mixing
these two up caused one of the two real bugs in \S\ref{sec:bugs}.

% ===========================================================================
\section{Code Walkthrough}
\label{sec:codewalkthrough}

\subsection{\texttt{src/lib/content.ts} --- the Data Layer}

\begin{lstlisting}
function loadYamlList<T>(filename: string): T[] {
  const filePath = path.join(CONTENT_DIR, filename);
  const raw = fs.readFileSync(filePath, 'utf8');
  const data = yaml.load(raw);
  if (!Array.isArray(data)) {
    throw new Error(`${filename} must parse to a YAML list at the top level`);
  }
  return data as T[];
}
\end{lstlisting}

Plain Node.js: read a file as text, parse it, sanity-check the shape ---
fail loudly at build time if someone breaks the YAML, rather than
silently rendering a blank page. \texttt{getExperience()} and
\texttt{getProjects()} wrap this with the specific filename and a
TypeScript type, so every page that imports them gets autocomplete and
type errors if a field is misspelled.

\texttt{getProjects()} also does the one piece of real logic in this
file: sorting --- featured projects first, then newest-by-date. That
logic lives here once, instead of being duplicated in every page that
lists projects (the same ``single source of truth'' principle from
\S\ref{sec:contentmodel}, applied to \emph{behavior} instead of
\emph{data}).

\texttt{formatMonthYear()} turns \texttt{"2026-07"} into \texttt{"Jul
2026"} for display, using the built-in \texttt{Intl}-backed
\texttt{toLocaleDateString}, so dates are stored in a
sortable/parseable machine format in the YAML but shown in a
human-readable format on the page.

\subsection{\texttt{src/components/} --- Layout, Nav, Footer, Cards}

\begin{itemize}[leftmargin=1.4em]
  \item \textbf{\texttt{Layout.astro}} is the page shell every page
  wraps itself in: emits \texttt{<html>/<head>} boilerplate once (meta
  tags, favicon, Open Graph tags for link previews, canonical URL for
  SEO), then renders \texttt{<Nav>}, \texttt{<slot />} (the page's own
  content), then \texttt{<Footer>}.
  \item \textbf{\texttt{Nav.astro}} builds its links from
  \texttt{import.meta.env.BASE\_URL} rather than hardcoding
  \texttt{/WebCV/...}, so the same code works whether the site is
  deployed at a subpath or a domain root. It marks the current page with
  \texttt{aria-current="page"} --- an accessibility attribute screen
  readers announce, which is also what the CSS uses to draw the orange
  underline.
  \item \textbf{\texttt{ExperienceCard.astro} / \texttt{ProjectCard.astro}}
  each take one entry as a prop and render it, with no knowledge of how
  many times they're used --- the pages loop and call them once per
  entry.
\end{itemize}

\subsection{\texttt{src/pages/} --- File-Based Routing and the Dynamic-Route Trick}

In Astro, the file path \emph{is} the URL --- no separate router config.
\texttt{src/pages/experience.astro} is served at \texttt{/experience}.

The interesting one is \file{src/pages/projects/[slug].astro} --- square
brackets mean ``dynamic route; the value in \texttt{[slug]} comes from
data, not a fixed filename.'' It exports \texttt{getStaticPaths()}:

\begin{lstlisting}
export function getStaticPaths() {
  return getProjects().map((project) => ({
    params: { slug: project.slug },
    props: { project },
  }));
}
\end{lstlisting}

At build time, Astro calls this once, gets back a list (currently 2
items --- \texttt{verifythevector} and \texttt{adversarial-robustness}),
and generates one actual \texttt{.html} file per item. Add a third
project to \texttt{projects.yaml} and the \emph{next build} produces a
third HTML file automatically --- no route to register, because the
route was never ``projects/verifythevector,'' it was always ``however
many projects exist.''

\subsection{\texttt{src/styles/global.css} --- Design Tokens, Dark Mode, Responsiveness}

The palette (a three-anime blend --- Naruto orange primary accent,
navy/cream base pair, sparing One Piece gold / Food Wars crimson
accents; exact colors and reasoning in \texttt{SPEC.md} \S9.1) is
implemented as CSS custom properties, not hardcoded hex values scattered
through the stylesheet:

\begin{lstlisting}
:root {
  --color-accent: #e8611c;
  --color-bg: #fbf6ee;
  --color-text: #232323;
  /* ...more tokens... */
}

@media (prefers-color-scheme: dark) {
  :root {
    --color-bg: #1b1f2a;
    --color-text: #edeae3;
    /* accent orange is NOT redefined -- stays the same in both modes */
  }
}
\end{lstlisting}

Every component then uses \texttt{var(--color-bg)} etc. instead of a
literal color --- a \textbf{design-token system}: define each
color/size once, by role (``background,'' ``accent,'' ``border'')
rather than by value, then reference the role everywhere. The payoff:
dark mode is a $\sim$10-line override block, not a second copy of every
color rule, because only the \emph{tokens} change; every rule using
\texttt{var(--color-bg)} picks up the new value automatically. This also
respects the visitor's OS-level light/dark preference automatically via
the \texttt{prefers-color-scheme} media query, with no toggle button or
JavaScript required.

Responsiveness follows the same ``define once'' instinct: layout uses
CSS Grid/Flexbox with relative sizing (e.g.
\texttt{repeat(auto-fill, minmax(280px, 1fr))} for card grids --- ``as
many 280px+ columns as fit,'' reflowing from multi-column on desktop to
single-column on phones with no media query needed) plus two
breakpoints (\texttt{480px}, \texttt{768px}) for the handful of things
that need explicit adjustment. \texttt{clamp()} is used for heading
sizes so text scales smoothly instead of jumping at breakpoints.

% ===========================================================================
\section{CI/CD: How a \texttt{git push} Becomes a Live Website}
\label{sec:cicd}

CI/CD (Continuous Integration / Continuous Deployment) means a pipeline
runs automatically and builds/publishes for you, consistently, every
time --- instead of manually running a build and manually uploading
files. Here that's a \textbf{GitHub Actions workflow}, defined in
\file{.github/workflows/deploy.yml}.

\begin{lstlisting}
on:
  push:
    branches: [main]     # runs automatically on every push to main
  workflow_dispatch:      # ...and can also be triggered manually

permissions:
  contents: read
  pages: write
  id-token: write         # OIDC auth to Pages, no manual secret needed

jobs:
  build:
    steps:
      - checkout            # clone the repo into the runner
      - setup-node          # install Node.js
      - npm ci               # install deps from the lockfile exactly
      - npm run build         # -> produces dist/
      - upload-pages-artifact # hand dist/ off to the next job

  deploy:
    needs: build            # doesn't start until build succeeds
    steps:
      - deploy-pages         # publish the artifact to GitHub Pages
\end{lstlisting}

Two jobs, not one, is a deliberate GitHub Pages convention:
\texttt{build} produces the static files and can run with narrow
permissions; \texttt{deploy} is the only step that needs permission to
publish, and it only runs \texttt{needs: build} --- if the build fails
(a broken YAML file, a typo in a component), deploy never runs and the
previously-published site stays live. Nothing broken ever gets
automatically published.

\begin{whybox}
\texttt{npm ci} (not \texttt{npm install}) matters here: it installs
\emph{exactly} the versions pinned in \texttt{package-lock.json}, so CI
uses the identical dependency versions as local development ---
reproducible builds, no ``works on my machine.'' The
\texttt{concurrency} block (\texttt{group: pages,
cancel-in-progress: true}) means if two pushes land close together, the
older in-flight deploy gets cancelled rather than both racing to
publish --- the site always ends up matching the latest commit.
\end{whybox}

GitHub Pages itself is a static file host built into GitHub --- free,
and pointed at ``GitHub Actions'' as its source (as opposed to ``a
specific branch's raw files,'' the older Pages model), which is what
lets this custom build step run at all.

% ===========================================================================
\section{Two Real Bugs Hit and Fixed}
\label{sec:bugs}

Small, concrete, and each shows a debugging \emph{process}, not just a
fix --- good material for ``tell me about a bug you debugged.''

\subsection{Bug 1: Broken Internal Links (\texttt{/WebCVexperience})}

\textbf{Symptom:} after the first successful build, every internal link
was missing a slash --- \texttt{href="}\file{/WebCVexperience}\texttt{"}
instead of \texttt{href="}\file{/WebCV/experience}\texttt{"}.

\textbf{Root cause:} \texttt{astro.config.mjs} had
\texttt{base: '/WebCV'} (no trailing slash). Astro exposes that value
unmodified as \texttt{import.meta.env.BASE\_URL}. Every link in the
codebase was built as \texttt{\`{}\$\{base\}experience\`{}} --- string
concatenation, expecting \texttt{base} to already end in \texttt{/}.
Without the trailing slash, that produced \texttt{"/WebCV" +
"experience"}.

\textbf{How it was found:} not by staring at code --- by checking the
actual build output. After \texttt{astro build}, every generated
\texttt{href="..."} was grepped and read. The build itself didn't error
(concatenating two strings is valid JS), so this was a ``silently wrong
output'' bug --- the kind that's easy to ship if you only check ``did
the build succeed'' instead of ``is the output actually correct.''

\textbf{Fix:} one character --- \texttt{base: '/WebCV/'} --- plus
rebuilding and re-grepping to confirm every link was correct, with a
comment left on that line explaining why the trailing slash is
required.

\subsection{Bug 2: Report/Slide PDF Links Pointed at Files That Don't Exist in the Deployed Site}

\textbf{Symptom:} \texttt{projects.yaml} originally had
{\footnotesize\texttt{report: "}\file{/content/Final\_Report\_VerifyTheVector.pdf}\texttt{"}}.
That path is valid \emph{in the repo}, but \texttt{content/} is never copied
into \texttt{dist/} (\S\ref{sec:repotour}) --- so that link would 404 on
the live site, even though nothing about the build failed.

\textbf{How it was found:} same instinct as Bug 1 --- after building,
what got written to \texttt{dist/} was actually inspected, tracing
whether every linked file existed there.

\textbf{Fix:} the three linked PDFs were copied into
\texttt{public/documents/} (the folder that \emph{does} get deployed)
and the YAML paths updated to point there, plus the page template
updated to prefix internal paths with the base path (same mechanism as
Bug 1), leaving external links (GitHub repo URLs) untouched.

\begin{whybox}[title={\textbf{The general lesson from both bugs}}]
A successful build (\texttt{exit code 0}) proves the code
\emph{compiled}, not that the output is \emph{correct}. Verifying this
project meant actually reading the generated HTML and following the
links, not just trusting a green checkmark.
\end{whybox}

% ===========================================================================
\section{Design-Process Decisions Worth Explaining}
\label{sec:designprocess}

\begin{itemize}[leftmargin=1.4em]
  \item \textbf{Why write \texttt{SPEC.md} before writing any code?} To
  separate ``what does this need to do'' from ``how do I build it,''
  and to have a written artifact of requirements (acceptance criteria,
  breakpoints, color palette) that can be checked against later instead
  of relying on memory --- standard practice on real engineering teams:
  a design doc precedes implementation so decisions get made
  deliberately, not accidentally while typing code.
  \item \textbf{Why \texttt{PLAN.md} in addition to \texttt{SPEC.md}?}
  \texttt{SPEC.md} answers ``what must be true when this is done'' and
  barely changes. \texttt{PLAN.md} answers ``what's actually done right
  now, and what did we decide along the way'' --- it changes constantly
  and is where decisions (``use Astro,'' ``use GitHub Pages,'' ``this
  repo is private, don't link it'') get written down \emph{once}.
  \item \textbf{Why per-project detail pages instead of just a grid?} A
  single grid page has to compress every project into a two-line
  summary. Once there's more than a couple of projects, a dedicated page
  per project lets each one carry its full writeup, tech stack, and
  links without cluttering the overview --- and gives each project its
  own shareable URL.
\end{itemize}

% ===========================================================================
\section{Glossary}
\label{sec:glossary}

\begin{longtable}{@{}p{4.0cm}p{10.5cm}@{}}
\toprule
\textbf{Term} & \textbf{Meaning} \\
\midrule
\endhead
Static site generation (SSG) &
  Building HTML files once, ahead of time, rather than per-request on a
  server (SSR) or in the browser via JavaScript (client-side
  rendering/SPA). \\[4pt]
Build time vs. runtime &
  ``Build time'' = when \texttt{astro build} runs, in CI, before anyone
  visits. ``Runtime'' = when a real visitor's browser loads the page.
  This project does essentially nothing at runtime beyond serving static
  files. \\[4pt]
CI/CD &
  Continuous Integration / Continuous Deployment --- automatically
  building and publishing on every change, instead of doing it by
  hand. \\[4pt]
YAML &
  A human-readable data format (like JSON, but supports comments and is
  less punctuation-heavy) --- used for \texttt{content/*.yaml}. \\[4pt]
Design tokens &
  Named variables for design values (colors, spacing) defined once and
  referenced everywhere, so changing the palette means editing one
  place. \\[4pt]
\texttt{prefers-color-scheme} &
  A CSS media query that detects the visitor's OS-level light/dark mode
  preference. \\[4pt]
OIDC &
  A way for a CI job to prove its identity to a service (here, GitHub
  Pages) without a manually-created, manually-rotated secret token. \\[4pt]
\texttt{getStaticPaths} &
  An Astro function that says ``generate one page per item in this
  list'' for a dynamic route --- the mechanism behind
  \texttt{/projects/[slug].astro} producing one real HTML file per
  project. \\[4pt]
Base path &
  The URL prefix a site is served under when it's not at a domain's
  root (here, \texttt{/WebCV/}, because GitHub Pages project sites live
  at \texttt{username.github.io/repo-name/}). \\
\bottomrule
\end{longtable}

% ===========================================================================
\section{Interview Cheat-Sheet}
\label{sec:cheatsheet}

\begin{description}[leftmargin=1.6em, style=nextline]
  \item[``Walk me through this project.''] Use the pitch in
  \S\ref{sec:pitch}, then let them steer with follow-ups.
  \item[``Why Astro instead of React/Next.js?''] The content barely
  changes and there's no interactivity that needs a client-side
  framework; Astro renders to plain HTML at build time so nothing ships
  to the browser except what's actually needed, while still giving
  reusable components and data-driven templating.
  \item[``How would you add a new project to the site?''] Copy the
  template block in \texttt{content/projects.yaml}, fill in the fields,
  save. Next build/deploy produces a new card on \texttt{/projects} and
  a whole new detail page automatically --- no code change.
  \item[``What happens when you push to \texttt{main}?''] Walk through
  \S\ref{sec:cicd}: a \texttt{build} job (\texttt{npm ci},
  \texttt{astro build}), then a \texttt{deploy} job that only runs if
  \texttt{build} succeeded, publishing to GitHub Pages. Live in under a
  minute.
  \item[``Tell me about a bug you ran into.''] Either bug in
  \S\ref{sec:bugs} --- both have a clear symptom, root cause, and fix,
  and both illustrate the same lesson: a successful build doesn't
  guarantee correct output.
  \item[``How is content separated from code?''] \S\ref{sec:contentmodel}:
  YAML files hold data, a loader reads and types it, page templates loop
  over whatever the loader returns.
  \item[``How does dark mode work?''] \S\ref{sec:codewalkthrough}: CSS
  custom properties define colors by role; a
  \texttt{prefers-color-scheme: dark} media query overrides just those
  variables; every rule referencing a variable updates automatically ---
  no JavaScript, no toggle.
  \item[``Why didn't you use a database / CMS?''] The content is small,
  changes rarely, and there's already a source of truth (plain text
  files in git). A database/CMS adds infrastructure and a network
  dependency to solve a problem version-controlled files already solve
  for free.
\end{description}

\end{document}
